
%%%%%%%%%%%%%%%%%%%%%%%%%%%%%%%%%%%%%%%%
%           main body
%%%%%%%%%%%%%%%%%%%%%%%%%%%%%%%%%%%%%%%%


%-----------------------------------
%	TITLE PAGE
%-----------------------------------

\title[Probability \& Statistics]{\LARGE \bfseries 第三部分\quad 概率与统计} % The short title appears at the bottom of every slide, the full title is only on the title page
\author[Conjugate Distributions] % Your institution as it will appear on the bottom of every slide, may be shorthand to save space
{\Large \bfseries \S 5. 共轭分布}
% \author{Singular Value Decomposition} % Your name
% \date{\today} % Date, can be changed to a custom date
\date{}

%------------------------------------------------

\begin{document}


%------------------------------------------------

\begin{frame}

\titlepage % Print the title page as the first slide

\end{frame}

% %------------------------------------------------

% \begin{frame}
% \frametitle{内容提要} % Table of contents slide, comment this block out to remove it
% \tableofcontents % Throughout your presentation, if you choose to use \section{} and \subsection{} commands, these will automatically be printed on this slide as an overview of your presentation
% \end{frame}

%-----------------------------------
%	PRESENTATION SLIDES
%-----------------------------------


%------------------------------------------------

\begin{frame}

这一章我们采用贝叶斯学派的记号以及术语。

\vspace{0.5em}

先来回忆一下贝叶斯公式
\[p(\theta \mid x)={\dfrac{p(x \mid \theta)\, p(\theta)}{ \int p(x \mid \theta')\,p(\theta')\,d\theta '}}\]
希望从数据$x$来推断参数$\theta$的分布（密度）。假设似然函数$p(x \mid \theta)$取定（这是一个合理的假设，因为一般地，由数据的产生过程，就能很好的确定似然函数）。

\end{frame}

%------------------------------------------------

\begin{frame}

\begin{block}{共轭分布，Conjugate Distribution}
如果后验分布$p(\theta \mid x)$与先验分布$p(\theta)$属于同类型的分布（参数一般不同），则
\begin{itemize}
\item 先验分布与后验分布被称为共轭分布
\item 先验分布被称为似然函数$p(x \mid \theta)$的共轭先验
\end{itemize}
\end{block}

\vspace{1em}
\pause

选取共轭先验的好处：
\begin{itemize}
\item 计算方便：分布形式确定，不用进行数值积分，方便迭代
\item 直观：迭代的时候，很容易看清楚似然函数如何更新先验分布
\end{itemize}

\end{frame}

%------------------------------------------------

\begin{frame}

\begin{block}{例子}
\begin{itemize}
\item Beta分布是二项似然分布的共轭先验
\end{itemize}

\vspace{1em}
\pause

事实上，取先验分布为Beta分布$Beta(\alpha,\beta)$, 即
\[p(\theta) = \dfrac{\theta^{\alpha-1} (1-\theta)^{\beta-1}}{\mathrm{B}(\alpha,\beta)}\]
似然分布为二项分布$B(n,\theta)$:
\[p(x \mid \theta) = C_n^x \theta^x (1-\theta)^x\]
\end{block}

\end{frame}

%------------------------------------------------

\begin{frame}

\begin{block}{例子（续）}
那么，后验分布有
\begin{align*}
p(\theta \mid x) \quad & {\color{red}\propto} \quad p(x \mid \theta) p(\theta) \\
& {\color{red}\propto} \quad C_n^x \theta^x (1-\theta)^{n-x} \dfrac{\theta^{\alpha-1} (1-\theta)^{\beta-1}}{\mathrm{B}(\alpha,\beta)} \\
& {\color{green}\propto} \quad \dfrac{\theta^{\alpha+x-1} (1-\theta)^{\beta+n-x-1}}{\mathrm{B}(\alpha+x,\beta+n-x)} \sim Beta(\alpha+x,\beta+n-x)
\end{align*}
${\color{red}\propto}$的系数为
\[\left( \int_0^1 C_n^x \dfrac{\theta'^{\alpha+x-1} (1-\theta')^{\beta+n-x-1}}{\mathrm{B}(\alpha,\beta)} d\theta' \right)^{-1} = \dfrac{\mathrm{B}(\alpha,\beta)}{C_n^x \mathrm{B}(\alpha+x,\beta+n-x)}\]
于是有${\color{green}\propto}$（实际是=）。
\end{block}

\end{frame}

%------------------------------------------------

\begin{frame}

\begin{block}{其他共轭先验}
\begin{center}
\renewcommand{\arraystretch}{1.3}
\begin{tabu} to 1\textwidth { | X[1.2,c] | X[0.8,c] | X[2,c] | }
\hline
先验分布 & 似然分布 & 后验分布 \\
$p(\theta)$ & $p(x \mid \theta)$ & $p(\theta \mid x)$ \\
\hline
$\operatorname{Beta}(\alpha,\beta)$ & $Ge(\theta)$ & $\operatorname{Beta}(\alpha+1,\beta+x-1)$ \\
\hline
$\operatorname{Beta}(\alpha,\beta)$ & $NB(r,\theta)$ & $\operatorname{Beta}(\alpha+x,\beta+r)$ \\
\hline
$\operatorname{Gamma}(\alpha,\beta)$ & $P(\theta)$ & $\operatorname{Gamma}(\alpha+x,\beta+1)$ \\
\hline
$\operatorname{Dir}(\boldsymbol\alpha)$ & $M(K,\mathbf{p})$ & $\operatorname{Dir}(\boldsymbol\alpha + \mathbf{x})$ \\
\hline
$N(\mu_{prior},\sigma^2_{prior})$ & $N(\theta,\sigma^2)$ & $N(\mu_{post},\sigma^2_{post})$ \\
\hline
\end{tabu}
\renewcommand{\arraystretch}{1}
\end{center}
最后一行参数的关系式为
\[\dfrac{\mu_{post}}{\sigma^2_{post}} = \dfrac{\mu_{prior}}{\sigma^2_{prior}} + \dfrac{x}{\sigma^2}, \quad \dfrac{1}{\sigma^2_{post}} = \dfrac{1}{\sigma^2_{prior}} + \dfrac{1}{\sigma^2}\]
\end{block}

\end{frame}

%------------------------------------------------

% \begin{frame}

% \begin{block}{Vandermonde行列式的计算}
% \begin{enumerate}
% % \addtocounter{enumi}{0}
% \item 把第$n-1$行乘以$-a_1$加到第$n$行，再把第$n-2$行乘以$-a_1$加到第$n-1$行。按照这种顺序依次向上，直到用第$2$行乘以$-a_1$加到第$1$行。由行列式性质，$V_n$的值不变。
% \[
% V_n = \det \begin{pmatrix}
% 1 & 1 & \cdots & 1 \\
% 0 & a_2-a_1 & \cdots & a_n-a_1 \\
% \vdots & \vdots & & \vdots \\
% 0 & a_2^{n-2}(a_2-a_1) & \cdots & a_n^{n-2}(a_n-a_1) \\
% \end{pmatrix}
% \]
% \end{enumerate}
% \end{block}

% \end{frame}

%------------------------------------------------
% % closing page
% \begin{frame}
% \HUGE{\centerline{\bfseries {\color{gray} The} {\color{zkblue} End}}}

% \pause

% \vspace{1.2em}

% \Large{\centerline{\bfseries {\color{gray}Thank} \quad {\color{zkblue} You}}}

% \end{frame}

%----------------------------------------------------------------------------------------

\end{document}
